% GENERATED FILE. Edit canonical lessons, then run npm run generate:books.
% canonical-source-hash: fnv1a64:678971957983365b
% canonical-lessons: BN-C344-rashtro, BN-C344-porichalona-kora, BN-C344-proshason, BN-C344-khomota, BN-C344-jonogon

\chapter{Who Governs}
\label{ch:bn-who-governs}
\hlchaptermodality{344}

\begin{chapteropening}
\textbf{By the end of this chapter:} \emph{I can say who holds power and who runs things.}

Complete the last lesson of chapter 344: I can say who holds power and who runs things.
\end{chapteropening}

\section[\texorpdfstring{\bn{রাষ্ট্র} (rāshṭrô)}{রাষ্ট্র (rāshṭrô)}]{\bn{রাষ্ট্র} (rāshṭrô) — a nation-state; the state}

\label{lesson:BN-C344-rashtro}



\begin{quote}
Before the new one: say the Bengali for to answer for, then the Bengali for to let off.
\end{quote}

\subsection*{You'll want to know: \bn{রাষ্ট্র}}
\textbf{\bn{রাষ্ট্র}} — \emph{rāshṭrô} — \textquotedblleft{}a nation-state; the state\textquotedblright{}.

\begin{quote}
\bn{রাষ্ট্রের} \bn{প্রধান} \bn{কাজ} \bn{নাগরিকদের} \bn{রক্ষা} \bn{করা।}
\end{quote}

\begin{quote}
\emph{rāshṭrer prôdhān kāj nāgôrikder rôkkhā kôrā.}
\end{quote}

\begin{quote}
The main job of the state is to protect its citizens.
\end{quote}

\begin{cousinweb}
\textbf{\bn{রাষ্ট্র}} (\emph{rāshṭrô}) is Sanskrit \emph{rāṣṭra}, a kingdom or realm (Monier-Williams). Modern Bengali uses it for the state as a political body, beside \bn{দেশ}, the country as a place.
\end{cousinweb}

\begin{grammarlens}[title={What you've built}]
One more word for giving reasons and opinions.
\end{grammarlens}

\subsection*{Your turn}
\begin{itemize}
\raggedright
  \item \emph{Say it:} \emph{rāshṭrô}
  \item \emph{Say it:} \emph{rāshṭrô}, once more
  \item \emph{Say it:} \emph{rāshṭrer prôdhān kāj nāgôrikder rôkkhā kôrā.} — The main job of the state is to protect its citizens.
  \item \emph{Recall:} say the Bengali for to answer for, then the Bengali for to let off, then say \emph{rāshṭrô} again
\end{itemize}

\begin{tcolorbox}[breakable,colback=teal!4,colframe=teal!35!black,title={Before you move on}]
What does \bn{রাষ্ট্র} mean? (\textquotedblleft{}A nation-state; the state\textquotedblright{}.) Say it once more.
\end{tcolorbox}
\section[\texorpdfstring{\bn{পরিচালনা} \bn{করা} (pôrichālônā kôrā)}{পরিচালনা করা (pôrichālônā kôrā)}]{\bn{পরিচালনা} \bn{করা} (pôrichālônā kôrā) — to run; to manage (an organisation)}

\label{lesson:BN-C344-porichalona-kora}



\begin{quote}
Before the new one: say the Bengali for to let off, then the Bengali for a nation-state.
\end{quote}

\subsection*{You'll want to know: \bn{পরিচালনা} \bn{করা}}
\textbf{\bn{পরিচালনা} \bn{করা}} — \emph{pôrichālônā kôrā} — \textquotedblleft{}to run; to manage (an organisation)\textquotedblright{}.

\begin{quote}
\bn{মীরা} \bn{একটা} \bn{ছোট} \bn{স্কুল} \bn{পরিচালনা} \bn{করে।}
\end{quote}

\begin{quote}
\emph{mīrā ekṭā chhoṭo skul pôrichālônā kôre.}
\end{quote}

\begin{quote}
Mira runs a small school.
\end{quote}

Past (he/she …): \bn{তিনি} \bn{কোম্পানিটা} \bn{পরিচালনা} \bn{করলেন}, \emph{tini kompāniṭā pôrichālônā kôrlen}.

\begin{cousinweb}
\textbf{\bn{পরিচালনা} \bn{করা}} (\emph{pôrichālônā kôrā}) is a learned formation, probably Bengali, from Sanskrit \emph{pari-} 'around' and \emph{cālana}, moving or setting going (Monier-Williams): to keep a whole body moving. \bn{চালানো}, to run a machine, is its everyday cousin.
\end{cousinweb}

\begin{grammarlens}[title={What you've built}]
One more word for giving reasons and opinions.
\end{grammarlens}

\subsection*{Your turn}
\begin{itemize}
\raggedright
  \item \emph{Say it:} \emph{pôrichālônā kôrā}
  \item \emph{Say it:} \emph{pôrichālônā kôrā}, once more
  \item \emph{Say it:} \emph{mīrā ekṭā chhoṭo skul pôrichālônā kôre.} — Mira runs a small school.
  \item \emph{Recall:} say the Bengali for to let off, then the Bengali for a nation-state, then say \emph{pôrichālônā kôrā} again
\end{itemize}

\begin{tcolorbox}[breakable,colback=teal!4,colframe=teal!35!black,title={Before you move on}]
What does \bn{পরিচালনা} \bn{করা} mean? (\textquotedblleft{}To run; to manage (an organisation)\textquotedblright{}.) Say it once more.
\end{tcolorbox}
\section[\texorpdfstring{\bn{প্রশাসন} (prôshāsôn)}{প্রশাসন (prôshāsôn)}]{\bn{প্রশাসন} (prôshāsôn) — an administration}

\label{lesson:BN-C344-proshason}



\begin{quote}
Before the new one: say the Bengali for a nation-state, then the Bengali for to run an organisation.
\end{quote}

\subsection*{You'll want to know: \bn{প্রশাসন}}
\textbf{\bn{প্রশাসন}} — \emph{prôshāsôn} — \textquotedblleft{}an administration\textquotedblright{}.

\begin{quote}
\bn{স্থানীয়} \bn{প্রশাসন} \bn{রাস্তাটা} \bn{সারাল।}
\end{quote}

\begin{quote}
\emph{sthānīyô prôshāsôn rāstāṭā sārālô.}
\end{quote}

\begin{quote}
The local administration repaired the road.
\end{quote}

\begin{cousinweb}
\textbf{\bn{প্রশাসন}} (\emph{prôshāsôn}) is Sanskrit \emph{praśāsana}, guidance or government (Monier-Williams). Bengali uses it for the officials who carry out the government's work.
\end{cousinweb}

\begin{grammarlens}[title={What you've built}]
One more word for giving reasons and opinions.
\end{grammarlens}

\subsection*{Your turn}
\begin{itemize}
\raggedright
  \item \emph{Say it:} \emph{prôshāsôn}
  \item \emph{Say it:} \emph{prôshāsôn}, once more
  \item \emph{Say it:} \emph{sthānīyô prôshāsôn rāstāṭā sārālô.} — The local administration repaired the road.
  \item \emph{Recall:} say the Bengali for a nation-state, then the Bengali for to run an organisation, then say \emph{prôshāsôn} again
\end{itemize}

\begin{tcolorbox}[breakable,colback=teal!4,colframe=teal!35!black,title={Before you move on}]
What does \bn{প্রশাসন} mean? (\textquotedblleft{}An administration\textquotedblright{}.) Say it once more.
\end{tcolorbox}
\section[\texorpdfstring{\bn{ক্ষমতা} (khômôtā)}{ক্ষমতা (khômôtā)}]{\bn{ক্ষমতা} (khômôtā) — power; capacity}

\label{lesson:BN-C344-khomota}



\begin{quote}
Before the new one: say the Bengali for to run an organisation, then the Bengali for an administration.
\end{quote}

\subsection*{You'll want to know: \bn{ক্ষমতা}}
\textbf{\bn{ক্ষমতা}} — \emph{khômôtā} — \textquotedblleft{}power; capacity\textquotedblright{}.

\begin{quote}
\bn{নির্বাচনের} \bn{পরে} \bn{নতুন} \bn{দল} \bn{ক্ষমতায়} \bn{এল।}
\end{quote}

\begin{quote}
\emph{nirbāchôner pôre nôtun dôl khômôtāy elô.}
\end{quote}

\begin{quote}
After the election a new party came to power.
\end{quote}

\begin{cousinweb}
\textbf{\bn{ক্ষমতা}} (\emph{khômôtā}) is Sanskrit \emph{kṣamatā}, ability or fitness (Monier-Williams). Bengali uses it for any capacity and for political power: \bn{ক্ষমতায়} \bn{আসা}, to come to power.
\end{cousinweb}

\begin{grammarlens}[title={What you've built}]
One more word for giving reasons and opinions.
\end{grammarlens}

\subsection*{Your turn}
\begin{itemize}
\raggedright
  \item \emph{Say it:} \emph{khômôtā}
  \item \emph{Say it:} \emph{khômôtā}, once more
  \item \emph{Say it:} \emph{nirbāchôner pôre nôtun dôl khômôtāy elô.} — After the election a new party came to power.
  \item \emph{Recall:} say the Bengali for to run an organisation, then the Bengali for an administration, then say \emph{khômôtā} again
\end{itemize}

\begin{tcolorbox}[breakable,colback=teal!4,colframe=teal!35!black,title={Before you move on}]
What does \bn{ক্ষমতা} mean? (\textquotedblleft{}Power; capacity\textquotedblright{}.) Say it once more.
\end{tcolorbox}
\section[\texorpdfstring{\bn{জনগণ} (jônôgôn)}{জনগণ (jônôgôn)}]{\bn{জনগণ} (jônôgôn) — the people; the public}

\label{lesson:BN-C344-jonogon}



\begin{quote}
Before the new one: say the Bengali for an administration, then the Bengali for power.
\end{quote}

\subsection*{You'll want to know: \bn{জনগণ}}
\textbf{\bn{জনগণ}} — \emph{jônôgôn} — \textquotedblleft{}the people; the public\textquotedblright{}.

\begin{quote}
\bn{সরকার} \bn{জনগণের} \bn{জন্য} \bn{কাজ} \bn{করে।}
\end{quote}

\begin{quote}
\emph{sôrkār jônôgôner jônyô kāj kôre.}
\end{quote}

\begin{quote}
The government works for the people.
\end{quote}

\begin{cousinweb}
\textbf{\bn{জনগণ}} (\emph{jônôgôn}) is Sanskrit \emph{jana}, people, with \emph{gaṇa}, a multitude (Monier-Williams): the whole body of the people, the public a government serves.
\end{cousinweb}

\begin{grammarlens}[title={What you've built}]
One more word for giving reasons and opinions.
\end{grammarlens}

\subsection*{Your turn}
\begin{itemize}
\raggedright
  \item \emph{Say it:} \emph{jônôgôn}
  \item \emph{Say it:} \emph{jônôgôn}, once more
  \item \emph{Say it:} \emph{sôrkār jônôgôner jônyô kāj kôre.} — The government works for the people.
  \item \emph{Recall:} say the Bengali for an administration, then the Bengali for power, then say \emph{jônôgôn} again
\end{itemize}

\begin{tcolorbox}[breakable,colback=teal!4,colframe=teal!35!black,title={Before you move on}]
What does \bn{জনগণ} mean? (\textquotedblleft{}The people; the public\textquotedblright{}.) Say it once more.
\end{tcolorbox}
